\section{Introduction}\label{sec:intro}
Block floating point represents a vector by integer coefficients and a common
power-of-two scale. The scale is usually selected from the block maximum,
so it depends on the same data that are subsequently rounded. This choice
couples the rounding errors. In a sum or a matrix--vector product, even a
small covariance between two coordinates can matter because it contributes
once for every coordinate pair. Predicting the variance of an individual
error is therefore insufficient for predicting the error of a reduction.

Classical statistical quantization models approximate residuals by centered
uniform noise \cite{widrow1996,widrow2008}. Block floating point roundoff
has a long history \cite{kalliojarvi1996}, and statistical noise models
continue to inform accelerator design \cite{song2018,lian2019}.
Microscaling formats add a contemporary motivation \cite{ocpmx2023}.
The shared-scale dependence is also
recognized: Rohani and Fakhraie \cite{rohani2006} discuss correlations in
dynamically rescaled structures, and Soloveychik et al.
\cite{soloveychik2022} use maximum order statistics to derive inner-product
error bounds and study block-size selection for Gaussian data. The latter
model approximates nonmaximum residuals by uniform noise. Fox and Lindstrom
\cite{fox2026zfp} analyze systematic bias in the ZFP compression pipeline.
These studies leave a finite-block question: how can one compute the error
covariance of a maximum-selected quantizer directly from the input law?

We derive such a calculation from one-dimensional truncated moments.
For independent, identically distributed inputs, conditioning on a binade
of the maximum excludes the event that every coordinate lies below its
lower endpoint. This exclusion gives an explicit nonpositive conditional
covariance. Variation of the conditional error mean across binades supplies
a competing nonnegative contribution. Their sum determines whether a
reduction is inflated or attenuated. The calculation also extends to
independent inputs with different distributions, for which the conditional
pair correction can have either sign. Conditioning on latent factors then
provides moments for correlated inputs and covariances between blocks.

The covariance matrix gives exact variance and mean-square error formulas
for deterministic weighted reductions. It also provides a useful
approximation guarantee: retaining the conditional moments while dropping
the negative pair correction overestimates the sum's variance inflation
factor by at most one. Thus the approximation becomes relatively accurate
when inflation is large, while the full formula resolves small inflation
and deflation. The complete covariance matrix supplies the moments needed
for propagation through a linear operator.

The large-block limit reveals a separate obstruction. Bounded support with
a nondegenerate limiting error gives unit inflation. With unbounded support,
the pair correlation tends to zero, yet the number of pairs can amplify it.
An exponential-input counterexample proves that inflation diverges along a
dyadic subsequence even though every input moment is finite. A decreasing
pair correlation therefore cannot justify an independent-error model for
large reductions.

Our numerical study checks the formulas by exact rational enumeration,
high-precision integration, and independently sampled blocks. Weighted
reductions expose the role of cancellation and bias; elliptic solves connect
the covariance to solution accuracy, block layout, and a fixed storage
budget. A train/validation experiment tests the same allocation after the
source law is estimated rather than supplied. Symmetric inputs and independently randomized stochastic rounding
provide controls through their zero pair covariance
\cite{connolly2021,croci2022}. These results complement probabilistic
analyses of arithmetic rounding \cite{higham2019,higham2020,hallman2023},
mixed-precision kernels \cite{fasi2023}, fixed-quantizer inner-product
distortion \cite{ang2026quant}, and state-conditioned reduction trees
\cite{sao2026trees}. The contribution here is the dependence created when
the input quantizer selects its own common scale.
